
\subsubsection{5.1 Pre-Validation Conditions (h-c1)}

NER F1 score: 0.96 (exceeds 0.90 threshold).

Wikipedia coverage: 1.0 (50/50 entities covered; exceeds 0.90 threshold).

Both measurement assumptions validated. Analysis proceeded to pattern detection.

\subsubsection{5.2 Attention Pattern Detection (h-e1)}

\textbf{Entity-error entropy:} mean = 0.062, median = 0.000, 75th percentile = 0.124 (N=50).

\textbf{Non-entity-error entropy:} mean = 0.300, median = 0.333, 75th percentile = 0.520 (N=23).

\textbf{Statistical test:} p = 7.53e-07 (two-sample t-test), Cohen's d = -1.13 (large effect size).

The null hypothesis (no difference) is rejected at p < 0.05. Entity-substitution errors exhibit significantly lower attention entropy than non-entity errors.

\textbf{Zero-entropy entity-errors:} 60\% of entity-error samples (30/50) exhibit exactly zero entropy, indicating deterministic attention away from the correct entity span.

\textbf{Sample loss:} 27\% of non-entity samples (27/50) were excluded due to unmappable entity spans (spaCy character-level spans to GPT-2 BPE token indices). The remaining 23 non-entity samples were sufficient for statistical power (p < 0.001).

\subsubsection{5.3 Classification Mechanism (h-m1)}

\textbf{Optimal threshold:} 0.32 (determined by grid search on training set, N=58).

\textbf{Training accuracy:} 81.0\%.

\textbf{Test accuracy:} 86.7\% (13/15 correct, N=15).

\textbf{Baseline accuracy:} 53.3\% (random classifier).

\textbf{Improvement:} +33.4 percentage points over baseline.

\textbf{Confusion matrix (test set):}

\begin{table}[htbp]
\centering
\resizebox{\columnwidth}{!}{%
\begin{tabular}{lll}
\toprule
 & Predicted Entity & Predicted Non-Entity \\
\midrule
Actual Entity & 10 & 0 \\
Actual Non-Entity & 2 & 3 \\
\bottomrule
\end{tabular}
}
\end{table}

\textbf{Precision (entity-error class):} 100\% (10/10).

\textbf{Recall (non-entity-error class):} 60\% (3/5).

The 86.7\% test accuracy exceeds the 70\% utility threshold, confirming that entropy-based classification is actionable for automated routing.

\subsubsection{5.4 Correction Routing (h-m2, Mock Results)}

\textbf{GPT-3.5 mock results:}

\begin{itemize}
\item Matched routing (RAG): 52\% success rate.
\item Mismatched routing (COT): 28\% success rate.
\item Difference: +24 percentage points.
\item Relative improvement: +85.7\%.
\end{itemize}

\textbf{Llama-2 mock results:}

\begin{itemize}
\item Matched routing (RAG): 42\% success rate.
\item Mismatched routing (COT): 22\% success rate.
\item Difference: +20 percentage points.
\item Relative improvement: +90.9\%.
\end{itemize}

Both models exceed the $\geq 20$ percentage point threshold. However, these results are synthetic (configurable success rates in ablation environment) and do not reflect real-world correction performance. No Wikipedia retrieval, LLM generation, or GPT-judge evaluation was deployed.

The experiment validates pipeline structure (dual-model framework, gate checking, matched versus mismatched comparison) but not correction effectiveness.
